\usepackage{xeCJK}
\usepackage{geometry}
\geometry{left=2cm,right=2cm,top=2.5cm,bottom=2.5cm}
\setlength{\headheight}{12.7pt}

\usepackage{titlesec}
\titleformat{\section}
  {\normalfont\large\bfseries}
  {\thesection}
  {1em}
  {}
\titlespacing*{\section}
  {0pt}
  {3.5ex plus 1ex minus .2ex}
  {2.3ex plus .2ex}

\usepackage{amsmath}
\usepackage{amssymb}
\usepackage{amsthm}
\usepackage{enumitem}
\setlist[enumerate]{itemsep=0pt,topsep=0pt,parsep=0pt,leftmargin=0pt,itemindent=2.5em}
\setlist[itemize]{itemsep=0pt,topsep=0pt,parsep=0pt,leftmargin=0pt,itemindent=2.5em}
\usepackage{fancyhdr}
\pagestyle{fancy}
\fancyhf{}
\usepackage{graphicx}
\usepackage{hyperref}
\usepackage{indentfirst}
\usepackage{lmodern}


\begin{document}
\setlength{\abovedisplayskip}{1pt}
\setlength{\belowdisplayskip}{1pt}

\section{拓扑空间}

% 保持乐观.

\hypertarget{sec:定义1.1}{}
\noindent\textbf{定义1.1 (拓扑空间)}

给定集合$X$和它的子集族$\mathcal{T}\subset\mathcal{P}(X)$. 如果:
\begin{enumerate}
    \item $\varnothing\in\mathcal{T},\,X\in\mathcal{T}$,
    \item 对$\mathcal{T}$中任意的$(U_i)_{i\in I}$,
    有$\displaystyle\bigcup_{i\in I}U_i\in\mathcal{T}$,
    \vspace{3pt}
    \item 对任意的$U_1,U_2\in\mathcal{T}$, 有$U_1\cap U_2\in\mathcal{T}$,
\end{enumerate}
就称$\mathcal{T}$是$X$上的一个\textbf{拓扑},
称$(X,\mathcal{T})$是一个\textbf{拓扑空间}.

在不产生歧义的情况下, 简称$X$是一个拓扑空间.

\vspace{10pt}

这时通常称$X$的元素为点, 称$X$的子集为点集.

\vspace{10pt}

\hypertarget{sec:定义1.2}{}
\noindent\textbf{定义1.2}

设$\mathcal{T}_1,\mathcal{T}_2$是$X$上的两个拓扑.
如果$\mathcal{T}_1\subset\mathcal{T}_2$, 就称$\mathcal{T}_1$比$\mathcal{T}_2$更粗,
或称$\mathcal{T}_2$比$\mathcal{T}_1$更细.

\vspace{10pt}

\hypertarget{sec:例1.1}{}
\noindent\textbf{例1.1 (平凡拓扑)}

给定集合$X$, $\{\varnothing,X\}$是$X$上最粗的拓扑, 称作\textbf{平凡拓扑}.

\vspace{10pt}

\hypertarget{sec:例1.2}{}
\noindent\textbf{例1.2 (离散拓扑)}

给定集合$X$, $\mathcal{P}(X)$是$X$上最细的拓扑, 称作\textbf{离散拓扑}.

\vspace{10pt}

\hypertarget{sec:例1.3}{}
\noindent\textbf{例1.3 (滤子拓扑)}

给定集合$X$上的\href{../../01 基础集合论/06 滤子/01 滤子#nameddest=sec:定义1.1}{滤子}
$F$, $F\cup\{\varnothing\}$是$X$上的拓扑, 称作\textbf{滤子拓扑}.





\vspace{20pt}

\section{开集}

\hypertarget{sec:定义2.1}{}
\noindent\textbf{定义2.1 (开集)}

给定拓扑空间$(X,\mathcal{T})$, 称$\mathcal{T}$中的点集为此空间的一个\textbf{开集},
称$\mathcal{T}$为此空间的\textbf{开集族}.

\vspace{10pt}

\hypertarget{sec:定理2.1}{}
\noindent\textbf{定理2.1 (开集的有限交封闭)}

给定拓扑空间$(X,\mathcal{T})$. 对任意有限个开集$(U_i)_{i<n}$,
它们的交$\displaystyle\bigcap_{i=0}^{n-1}U_i$也是开集.

\noindent{\kaishu 证明.}

首先, $n=1$的情况显然; 其次, 假设任取$(U_i)_{i<n}$
有$\displaystyle\bigcap_{i=0}^{n-1}U_i\in\mathcal{T}$, 那么任取$(U_i)_{i<n+1}$有:
\vspace{3pt}
\begin{equation*}
    \bigcap_{i=0}^nU_i=\left(\bigcap_{i=0}^{n-1}U_i\right)\cap U_n\in\mathcal{T}.
\tag*{\qedsymbol}\end{equation*}





\newpage

\section{闭集}

\hypertarget{sec:定义3.1}{}
\noindent\textbf{定义3.1 (闭集)}

给定拓扑空间$(X,\mathcal{T})$.

对任意点集$V$, 如果$V^c\in\mathcal{T}$, 就称$V$是此空间的一个\textbf{闭集},
称闭集的全体为此空间的\textbf{闭集族}.

\vspace{10pt}

可以根据闭集族来定义一个拓扑.

\vspace{10pt}

\hypertarget{sec:定义3.2}{}
\noindent\textbf{定义3.2 (闭集公理)}

给定集合$X$和它的子集族$\mathcal{C}\subset\mathcal{P}(X)$.
称$\mathcal{C}$满足\textbf{闭集公理}, 如果:
\begin{enumerate}
    \item $\varnothing\in\mathcal{C},\,X\in\mathcal{C}$,
    \item 对$\mathcal{C}$中任意的$(V_i)_{i\in I}$,
    有$\displaystyle\bigcap_{i\in I}V_i\in\mathcal{C}$,
    \vspace{3pt}
    \item 对任意的$V_1,V_2\in\mathcal{C}$, 有$V_1\cup V_2\in\mathcal{C}$.
\end{enumerate}

\vspace{10pt}

\hypertarget{sec:定理3.1}{}
\noindent\textbf{定理3.1}

给定集合$X$和它的子集族$\mathcal{C}\subset\mathcal{P}(X)$, 那么:

$\mathcal{C}$满足闭集公理当且仅当$X$上存在唯一的拓扑$\mathcal{T}$
使得$\mathcal{C}$是$(X,\mathcal{T})$的闭集族.

\noindent{\kaishu 证明.}

\noindent{\kaishu 必要性.}

假设$\mathcal{C}$满足闭集公理, 定义$\mathcal{T}=\{V^c\mid V\in\mathcal{C}\}$. 易证:
\begin{enumerate}
    \item $\varnothing=X^c\in\mathcal{T},\,X=\varnothing^c\in\mathcal{T}$,
    \vspace{3pt}
    \item 对$\mathcal{T}$中任意的$(U_i)_{i\in I}$,
    $\displaystyle\left(\bigcup_{i\in I}U_i\right)^c
    =\bigcap_{i\in I}U_i^c\in\mathcal{C}$, 所以
    $\displaystyle\bigcup_{i\in I}U_i\in\mathcal{T}$,
    \vspace{3pt}
    \item 对任意的$U_1,U_2\in\mathcal{T}$,
    有$(U_1\cap U_2)^c=(U_1^c\cup U_2^c)\in\mathcal{C}$, 所以
    $U_1\cap U_2\in\mathcal{T}$.
\end{enumerate}
即$\mathcal{T}$是$X$上的拓扑.
进一步, 显然$(X,\mathcal{T})$的闭集恰好是$\mathcal{C}$中的点集,
即$\mathcal{C}$是$(X,\mathcal{T})$的闭集族.

如果拓扑$\mathcal{T}_1,\mathcal{T}_2$的空间的闭集族都是$\mathcal{C}$,
那么对任意的$U\in\mathcal{T}_1$有
\[
    U^c\in\mathcal{C}\quad\implies\quad U\in\mathcal{T}_2,
\]
即$\mathcal{T}_1\subset\mathcal{T}_2$, 反之同理. 于是这样的拓扑唯一.

\noindent{\kaishu 充分性.}

假设$\mathcal{C}$是拓扑空间$(X,\mathcal{T})$的闭集族, 则:
\begin{enumerate}
    \item $\varnothing=X^c\in\mathcal{C},\,X=\varnothing^c\in\mathcal{C}$,
    \vspace{3pt}
    \item 对$\mathcal{C}$中任意的$(V_i)_{i\in I}$,
    $\displaystyle\left(\bigcap_{i\in I}V_i\right)^c
    =\bigcup_{i\in I}V_i^c\in\mathcal{T}$, 所以
    $\displaystyle\bigcap_{i\in I}V_i\in\mathcal{C}$,
    \vspace{3pt}
    \item 对任意的$V_1,V_2\in\mathcal{C}$,
    有$(V_1\cup V_2)^c=(V_1^c\cap V_2^c)\in\mathcal{T}$, 所以
    $V_1\cup V_2\in\mathcal{C}$. \hfill $\qedsymbol$
\end{enumerate}

\vspace{10pt}

\hypertarget{sec:定理3.2}{}
\noindent\textbf{定理3.2 (闭集的有限并封闭)}

给定拓扑空间$(X,\mathcal{T})$. 对任意有限个闭集$(V_i)_{i<n}$,
它们的并$\displaystyle\bigcup_{i=0}^{n-1}V_i$也是闭集.

\noindent{\kaishu 证明.}

每个$V_i^c$都是开集, 所以$\displaystyle\left(\bigcup_{i=0}^{n-1}V_i\right)^c
=\bigcap_{i=0}^{n-1}V_i^c$也是开集,
即$\displaystyle\bigcup_{i=0}^{n-1}V_i$是闭集. \hfill $\qedsymbol$

% 焦虑, 疲惫, 感伤.
% 2026.05.25

% 写完了毕设论文的第一稿, 最近最轻松的一晚. 希望毕业顺利.
% 有点想念数学了.
% 2026.05.27

% 论文似乎改完了?
% 事已至此, 再休息一会儿吧.
% 2026.05.31

\end{document}